\begin{abstract}
Multimodal learning commonly uses partial information decomposition (PID) synergy to quantify how much a task requires
cross-modal fusion. For deterministic answers, however, PID depends only on per-modality answer statistics, so it
cannot see whether answers require a non-additive interaction. For tasks with a known answer for every input pair, we
use a model-free split of the answer table into an additive rule, which scores each modality separately, and
exceptions, which no additive rule can answer. Computed exactly by integer programming, the exception share is the
minimum error of any additive model and, by a classical result, is zero exactly when the table has no XOR square. It
also sets a floor on how much an accurate model must lose when its cross-modal interactions are removed. On 99 fully
labelled answer tables from CLEVR, GQA, CLOVE and POPE, four 7--8B VLMs lose accuracy without interactions roughly in
proportion to this requirement: an accuracy-matched null scorer reproduces most of the relation, and only one model
clearly relies on interactions beyond it. In twelve continual-learning methods, exceptions on compositional CLEVR tasks
are forgotten faster and learned worse than rule cells with the same answer and the same difficulty, a loss that
standard backward transfer does not show.
\end{abstract}
